\uuid{av7Q}
\titre{ Détection d'une erreur dans une boucle d'entraînement }
\niveau{L3}
\module{Analyse de données}
\chapitre{Réseaux de neurones}
\sousChapitre{Descente de gradient, implémentation}
\theme{Montée de gradient, signe de la mise à jour, divergence}
\auteur{Maxime Nguyen}
\datecreate{2026-03-19}
\organisation{AMSCC}
\difficulte{2}

\contenu{
\texte{
  On considère la boucle d'entraînement suivante pour une régression linéaire :
\begin{verbatim}
for it in range(num_iter):
    y_hat = X @ w + b
    diff  = y_hat - y
    cost  = (1/(2*m)) * np.sum(diff**2)

    dw = (1/m) * X.T @ diff
    db = (1/m) * np.sum(diff)

    w += lr * dw   # <- ?
    b -= lr * db
\end{verbatim}
}
\begin{enumerate}

\item
  \question{Ce code contient une erreur. Laquelle, et quel effet a-t-elle sur l'entraînement ?}
  \reponse{
    La ligne \texttt{w += lr * dw} effectue une \emph{montée} de gradient sur $w$ au lieu d'une descente.
    La mise à jour correcte est \texttt{w -= lr * dw}.

    En conséquence, $w$ est mis à jour dans la direction qui \emph{augmente} le coût $J$,
    ce qui provoque la divergence de l'entraînement : le coût croît à chaque itération au lieu de décroître.
    Seul $b$ est mis à jour correctement (\texttt{b -= lr * db}).
  }

\end{enumerate}
}
